Multimodalita znamená zapojení více modalit (text, obraz, rukopis, interaktivní simulace) do zadání i sběru důkazů; umožňuje požadovat artefakty procesu jako doplněk k textovým odpovědím a tím zvyšuje validitu hodnocení \cite{kb_ai_101_tipu_pdf_04e4a115_page_8_1}.
Praktické postupy pro učitele: \begin{itemize}
\item Rukopis→OneNote: studenti píší rukou, pak použijte \textit{Kreslení}→\textit{Rukopis na text} pro strojově čitelný záznam \cite{kb_ai_101_tipu_pdf_04e4a115_page_13_2}.
\item Interaktivní světy (Minecraft Education): použijte připravené světy/lekce z Hour of Code a reflektujte učení diskusí \cite{kb_ai_101_tipu_pdf_04e4a115_page_15_2}.
\item Kombinovaná zadání: požadujte text + screenshot/export z OneNote jako artefakt procesu pro ověření autenticity \cite{kb_ai_101_tipu_pdf_04e4a115_page_8_1,kb_ai_101_tipu_pdf_04e4a115_page_13_2}.
\end{itemize}
Krátký checklist pro návrh multimodálních úloh: \begin{enumerate}
\item Určete modality; \item Dejte technické instrukce (export OneNote, screenshot, nahrání Minecraft artefaktů) \cite{kb_ai_101_tipu_pdf_04e4a115_page_13_2,kb_ai_101_tipu_pdf_04e4a115_page_15_2}; \item Požadujte artefakt procesu + popis workflow \cite{kb_ai_101_tipu_pdf_04e4a115_page_8_1}.
\end{enumerate}